\documentclass[11pt,a4paper]{article}
\usepackage[T1]{fontenc}
\usepackage[utf8]{inputenc}
\usepackage{lmodern}
\usepackage{amsmath,amssymb}
\usepackage[margin=25mm]{geometry}
\usepackage[hidelinks]{hyperref}
\hypersetup{pdftitle={Leibniz's law of continuity, read on records, selects a positive action floor for motions},pdfauthor={navstokgap research working notes}}
\newcommand{\tightlist}{\setlength{\itemsep}{0pt}\setlength{\parskip}{0pt}}
\setlength{\emergencystretch}{3em}
\setcounter{secnumdepth}{0}
\begin{document}
\hypertarget{leibnizs-law-of-continuity-read-on-records-selects-a-positive-action-floor-for-motions}{%
\section{Leibniz's law of continuity, read on records, selects a
positive action floor for
motions}\label{leibnizs-law-of-continuity-read-on-records-selects-a-positive-action-floor-for-motions}}

\textbf{Result, 2026-09-29 (Claude; written derivation and passage
readings, unrefereed).} In the mark model of the
\href{planck-gap-paper.md}{Planck paper} (§3), the least error with
which any admissible protocol decides Newton's comparison, inertia
against a constant force \(F\) over a cell of duration \(\tau\) with
unknown initial position and velocity, is

\[P_*(F)=\Phi\Bigl(-\tfrac12\sqrt{K_\tau/\kappa}\Bigr),\qquad K_\tau=\frac{F^2\tau^3}{24\,m},\]

for a mark floor \(\kappa>0\), while \(P_*(F)=0\) for every \(F\ne0\)
when \(\kappa=0\); at \(F=0\) the cases coincide and \(P_*=\frac12\).
\textbf{The best verdict is continuous in the force at \(F=0\) if and
only if \(\kappa>0\).} The same holds on Zeno's rung, rest against
uniform motion, with \(mv^2\tau/(8\kappa)\) in place of
\(K_\tau/(4\kappa)\). Static shapes behave differently: there the best
verdict jumps at zero for every \(\kappa\), because reading a shape
disturbs nothing.

Leibniz stated the law of continuity in July 1687, the year of the
\emph{Principia}, as a test of laws of motion: when two cases approach
and are lost in one another, their outcomes must do the same, and he
describes the nearly equal case of a collision as one ``dont à peine ce
cas peut estre distingué'' (§2). Read with outcomes measured by how well
they can be told apart, the law holds for Newton's comparison exactly
when \(\kappa>0\). Read with outcomes measured by positions and
velocities, it holds on both branches. The fork that every route of this
programme reaches (positivity as a premise, the zero branch admissible)
is therefore located inside one Newton-age text, and the programme's
thesis, that quantization is a consistency condition on laws which
Newton's continuum limit missed, has a period form: Leibniz's ``preuve
ou examen'' of laws of motion, applied to their records. His
\emph{petites perceptions} (1704) supply the graded, additive
discernibility that the records reading needs (§3). The Newton-age
dispute over indivisibles, Galileo, Cavalieri and Guldin between 1621
and 1647, concerned static figures, where the theorem imposes no floor
and Newton's limit doctrine is the right answer (§4).

\hypertarget{the-proposition}{%
\subsection{1. The proposition}\label{the-proposition}}

\textbf{Setting.} The statistical model of the Planck paper, §3: mark
\(j\) at time \(t_j\) returns the position with Gaussian error of
standard deviation \(\delta_j\) and delivers an independent Gaussian
impulse of standard deviation \(\Delta_j\); every available mark obeys
\(\delta_j\Delta_j\ge\kappa\); protocols are non-adaptive and tests are
invariant under the unknown \(y_0,v_0\). The optimal invariant test errs
with probability \(\Phi(-d/2)\), and Theorem 2 of the paper gives
\(d^2\le K_\tau/\kappa\) for every protocol, with a sharp constant.

\textbf{Proposition L.} Let \(P_*(F)\) be the infimum of the error over
admissible protocols and invariant tests. Then
\(P_*(F)=\Phi(-\frac12\sqrt{K_\tau/\kappa})\) if \(\kappa>0\), and
\(P_*(F)=0\) for \(F\ne0\) if \(\kappa=0\). For Zeno's rung (rest
against speed \(v\), unknown initial position),
\(P_*(v)=\Phi(-\sqrt{mv^2\tau/(8\kappa)})\) if \(\kappa>0\) and
\(P_*(v)=0\) for \(v\ne0\) if \(\kappa=0\). For a static taper
\(\vartheta\) read at two heights, \(P_*(\vartheta)=0\) for
\(\vartheta\ne0\) and every \(\kappa\ge0\).

\emph{Proof.} For \(\kappa>0\), \(\Phi(-d/2)\) decreases in \(d\), and
the supremum of \(d^2\) over protocols is \(K_\tau/\kappa\) by the
sharpness in Theorem 2; hence the infimum of the error. For
\(\kappa=0\), marks with \(\Delta_j=0\) and arbitrary \(\delta_j>0\) are
admissible; take three marks at distinct times and weights \(u\) with
\(\sum u_i=\sum u_it_i=0\), so that
\(u^{\mathsf T}P=\frac F{2m}\sum u_it_i^2\ne0\) while
\(u^{\mathsf T}\Sigma u=\sum\delta_i^2u_i^2\to0\); then \(d\to\infty\).
Zeno's rung is Theorem 9(i) of the paper, whose bound
\(d^2\le mv^2\tau/(2\kappa)\) two marks attain. The static case is
Theorem 9(iii): with \(N\) marks of resolution \(\delta\) at each
height, \(d^2=N\vartheta^2\Delta h^2/(2\delta^2)\), unbounded in \(N\)
whatever \(\kappa\) is. \(\square\)

For \(\kappa>0\) the verdict is Lipschitz at \(F=0\) with constant
\(\tau^{3/2}/(2\sqrt{2\pi}\sqrt{24m\kappa})\), which diverges as
\(\kappa\to0\): the zero branch is the singular limit of a family of
continuous verdict functions, the pattern of an order of limits. Quantum
mechanics agrees with the model on the split between the two kinds of
case. A static parameter can be estimated without limit by repetition,
while the dynamical comparison in a fixed window has the
\(N\)-independent bound that the paper's quantum results supply.

\hypertarget{leibniz-1687-the-law-of-continuity-as-a-test-of-laws}{%
\subsection{2. Leibniz, 1687: the law of continuity as a test of
laws}\label{leibniz-1687-the-law-of-continuity-as-a-test-of-laws}}

Leibniz, ``Lettre de M. L. sur un principe general utile à l'explication
des loix de la nature'', \emph{Nouvelles de la République des Lettres},
July 1687, in Gerhardt, \emph{Die philosophischen Schriften} III (1887),
51--55
(\href{../docs/classics/Leibniz_PrincipeGeneral_1687_Gerhardt3_OCR.md}{companion};
passage, OCR normalized by hand).

\begin{itemize}
\tightlist
\item
  \textbf{The text.} ``Lorsque la difference de deux cas peut estre
  diminuée au dessous de toute grandeur donnée \emph{in datis} ou dans
  ce qui est posé, il faut qu'elle se puisse trouver aussi diminuée au
  dessous de toute grandeur donnée \emph{in quaesitis} ou dans ce qui en
  resulte \ldots{} \emph{Datis ordinatis etiam quaesita sunt ordinata}''
  (p.~52). ``Le repos peut estre consideré comme une vistesse infinement
  petite, ou comme une tardité infinie'', so that the rule of rest must
  be a particular case of the rule of motion, ``autrement \ldots{} ce
  sera une marque asseurée, que les regles sont mal concertées''
  (pp.~52--53). Against Descartes's second rule of collision, a body B
  made ``aussi petite que l'on voudra'' larger than C should reflect a
  little less, and C a little more, ``qu'au cas de l'égalité dont à
  peine ce cas peut estre distingué'' (p.~53). And the exemption: ``dans
  les choses composées quelques fois un petit changement peut faire un
  grand effect, comme par exemple une estincelle tombant dans une grande
  masse de la poudre à canon'', while ``à l'égard des principes ou
  choses simples, rien de semblable ne sçauroit arriver'' (p.~54).
\item
  \textbf{What it commits him to.} A law of motion is tested from
  outside, before any inner discussion (``preuve ou examen \ldots{}
  avant même que de venir à une discussion interieure''), by whether its
  outcomes converge when its cases do. Rest is the limit of slow motion
  and equality the limit of inequality, and a rule that jumps at either
  is ill made. Composite amplifiers may magnify small differences into
  large effects; simple principles may not jump.
\item
  \textbf{What the theorem does with it.} Leibniz applies the test to
  collision velocities, the geometric reading, and Newton's laws pass it
  on both branches. His own phrase for the nearly equal case, however,
  is a statement about discernibility, and Proposition L shows what the
  test demands when outcomes are measured that way: for Newton's
  comparison and for Zeno's rung, continuity holds exactly when
  \(\kappa>0\). The exemption marks the boundary of the argument. An
  instrument is a composite amplifier, and finite amplification is
  allowed; the zero branch needs the supremum over all instruments the
  laws permit to jump, a jump in what the laws allow to be recorded,
  which is the level of principle if the mark trade-off is a law.
  Whether the verdict of the best possible record counts among ``ce qui
  en resulte'' is the premise, now stated in Leibniz's terms. The
  theorem also declines one extension: for static shapes the records
  reading fails on both branches, and Leibniz's letter applies the law
  to figures geometrically (the ellipse passing into the parabola,
  p.~52), which Proposition L shows to be the only reading that holds
  there.
\end{itemize}

Prior art, at metadata level: Mortensen, ``The Leibniz Continuity
Condition, Inconsistency and Quantum Dynamics'',
\href{https://doi.org/10.1023/A:1004275928191}{\emph{J. Philos. Logic}
\textbf{26} (1997), 377--389}, treats the condition against quantum
jumps at the instant of change; the pairing of continuity with the
identity of indiscernibles is the subject of Russell's chapter ``The
Identity of Indiscernibles and the Law of Continuity'' in \emph{A
Critical Exposition of the Philosophy of Leibniz}. The operational
reading used here is the base-norm topology on states, in which
classical pure states are pairwise perfectly distinguishable; compare
the \href{quantum-exclusion-premises.md}{Hardy audit}, whose continuity
axiom concerns reversible paths between pure states.

\hypertarget{leibniz-1704-small-perceptions-add}{%
\subsection{3. Leibniz, 1704: small perceptions
add}\label{leibniz-1704-small-perceptions-add}}

Leibniz, \emph{Nouveaux essais}, Préface, written 1703--1705, in
Gerhardt V (1882), 47--49
(\href{../docs/classics/Leibniz_NouveauxEssais_Preface_Gerhardt5_OCR.md}{companion};
passage, OCR normalized by hand).

\begin{itemize}
\tightlist
\item
  \textbf{The text.} To hear the roar of the sea ``il faut bien qu'on
  entende les parties qui composent ce tout, c'est à dire les bruits de
  chaque vague \ldots{} Car il faut qu'on en soit affecté un peu par le
  mouvement de cette vague \ldots{} autrement on n'auroit pas celle de
  cent mille vagues, puisque cent mille riens ne sauroient faire quelque
  chose'' (p.~47). ``La nature ne fait jamais des sauts \ldots{} jamais
  un mouvement ne naist immediatement du repos ny s'y reduit que par un
  mouvement plus petit'' (p.~49). ``En vertu des variations insensibles,
  deux choses individuelles ne sauroient estre parfaitement semblables''
  (p.~49).
\item
  \textbf{What it commits him to.} Every effect, however small,
  registers a little; perception is graded and additive; noticeable
  perceptions come by degrees from insensible ones; distinct individuals
  always differ, by differences that can be insensible.
\item
  \textbf{What the theorem does with it.} The cut measure makes the
  addition exact. Newton's cell action \(K_\tau\) is spent additively
  over any sequence of cuts, a share \(3s(1-s)K\) per cut at fraction
  \(s\) (\href{cut-measure-newton.md}{cut-measure note}, Theorems 1--2),
  and it is the Cameron--Martin energy of the force's shift (Proposition
  4), so each cut carries a finite share of the whole comparison's
  discernibility and the shares add to \(K_\tau/\kappa\). With
  \(\kappa>0\) each cut is a \emph{petite perception} of the force,
  insensible alone below the mark mesh and effective in the assembly;
  with \(\kappa=0\) a single cut decides, and records of motion have no
  insensible parts. The zero branch can still place insensibility in the
  perceiver, as Leibniz's confusion of created substances does, graded
  from one substance to another. The
  \href{stochastic-route-velocitas-ultima.md}{stochastic route} shows
  how composition turns graded constants into one (\(\kappa=mD\) for
  every body), so the step from ``every created perceiver is confused''
  to one universal floor is available as a reading; it is recorded here
  as a reading. The maxim on motion from rest is Galileo's passage
  through all degrees of slowness, stated at the beginning of motion,
  the regime of Lemma X (``ipso motus initio''); Zeno's rung of
  Proposition L is its records version.
\end{itemize}

\hypertarget{the-newton-age-dispute-over-indivisibles-16211647}{%
\subsection{4. The Newton-age dispute over indivisibles,
1621--1647}\label{the-newton-age-dispute-over-indivisibles-16211647}}

Book I's closing scholium prefers limits to indivisibles, ``quoniam
durior est indivisibilium Hypothesis'', without naming the dispute it
settles. The primary texts are held in \texttt{docs/classics}.

\begin{itemize}
\tightlist
\item
  \textbf{Galileo's bowl and cone} (\emph{Discorsi}, Giornata prima,
  Favaro VIII, 74--76, 78--80;
  \href{../docs/classics/Galileo_Discorsi_GiornataPrima_Favaro_it_wikisource.md}{companion}).
  Every horizontal section of the bowl (cylinder minus hemisphere)
  equals the section of the inscribed cone, so the surfaces are ``sempre
  eguali, e \ldots{} diminuendosi sempre egualmente, vadano a terminare
  l'una in un sol punto e l'altra nella circonferenza d'un cerchio
  \ldots{} perché in questa consequenza sola versa la nostra
  maraviglia''. His conclusion: ``questi attributi di maggioranza,
  minorità ed egualità non convenghino a gl'infiniti''.
  \emph{Commitment:} comparison fails among infinites and indivisibles.
  \emph{The theorem:} the bowl is Democritus's third (bowl and cone are
  each a third of the cylinder), a static figure, so Theorem 9(iii) and
  Proposition L put no floor on it, and Newton's sentence that ultimate
  ratios are limits and never ratios of ultimate quantities dissolves
  the wonder: the ratio of sections is 1 at every cut, and point and
  circle are never compared. Leibniz reverses Galileo's conclusion in
  1687: equality is an infinitely small inequality, so the attributes
  extend by continuity.
\item
  \textbf{Cavalieri's reply} (letter 2992 of 1634, Opere XVI 136--138;
  \href{../docs/classics/Galileo_Opere_XVI_indivisibles_letters_1634-1636_OCR.md}{companion}).
  Equal parts are removed ``essendo noi arrivati al nullo piano tanto
  nel cono quanto nella scodella'', and ``non mi dichiaro di componere
  il continuo d'indivisibili, ma solo mostro che i continui hanno la
  proportione delli aggregati di questi indivisibili''.
  \emph{Commitment:} ratios of aggregates of sections, with no claim of
  composition. \emph{The theorem:} the cut measure keeps exactly this
  stance, a statement about aggregates of cuts (shares summing to
  \(K_\tau\) for every schedule) that composes nothing, and the mark
  mesh bounds what can be exhibited without bounding what can be cut.
\item
  \textbf{Galileo's wheel} (Favaro VIII, 68--72, same companion). A
  smaller concentric wheel covers the larger wheel's line with its own
  sides ``con l'interposizione di cento mila spazii vacui traposti'',
  and for circles ``sì come i lati non son quanti, ma bene infiniti,
  così gl'interposti vacui non son quanti, ma infiniti''; Simplicio
  hears in it ``quei vacui disseminati di certo filosofo antico'', and
  Salviati's retort about the denier of Providence names Epicurus.
  \emph{Commitment:} a continuum composed of infinitely many
  unquantified indivisibles, partly full and partly void, crossed by
  leaps across the voids. \emph{The theorem:} Galileo publishes
  al-Nazzam's leap in 1638 as consistent. The paper's reading of the
  leap applies: below the mesh a record is consistent with a leap, the
  interval stays divisible, and the crossing is exhibited in finitely
  many steps.
\item
  \textbf{Guldin} (\emph{Centrobaryca} IV, 1641, preface p.~4;
  \href{../docs/classics/Guldin_Centrobaryca_LiberIV_1641_pages.md}{companion}):
  ``Galileus profecto in eodem Dialogo de Motu locali, disputans de
  infinito, de proprietatibus finitorum, quas infinitis applicare minime
  liceat, contra ipsum concludit.'' \emph{Commitment:} Galileo's own
  principle refutes Cavalieri. \emph{The theorem:} it sides with Guldin
  and Newton on composition (no least magnitude follows from \(\kappa\))
  and with Cavalieri on ratios of aggregates.
\end{itemize}

The four entries share one feature that sharpens the paper's §9. The
seventeenth-century dispute Newton answered concerned static figures,
and there the theorem agrees with his answer and imposes nothing. The
floor enters only where a record carries back-action, in the lemmas
applied to motion, and the dispute about motion is the one he passed
over.

\hypertarget{consequence-for-state}{%
\subsection{5. Consequence for STATE}\label{consequence-for-state}}

For the Newton goal: positivity of the action floor is equivalent, in
the mark model, to Leibniz's 1687 law of continuity for laws of motion
with outcomes measured by discernibility (Proposition L). The zero
branch keeps the geometric reading, which Leibniz himself applied, so
the fork stands; it is now a choice between two readings of one
Newton-age principle, with Leibniz's phrase ``à peine distingué'' and
his doctrine of small perceptions on the side of the floor, and statics
showing why the reading must be confined to motions. For the scholion:
two Leibniz entries and the Galileo--Cavalieri--Guldin layer are
supplied with their three obligations; a pointer belongs in the paper's
§9. To be refereed (Astra or Fable): the proposition, the passage
readings, and whether the records reading is fair to Leibniz.
\end{document}
